\documentclass[a4paper]{article}
\usepackage[utf8]{inputenc}
\usepackage{graphicx}
\usepackage{blindtext}
\usepackage{hyperref}
\usepackage{float}
\usepackage{caption}
\usepackage{subcaption}
\usepackage{indentfirst}
\usepackage{setspace}
\usepackage{pdfpages}
\usepackage{svg}
\usepackage{lipsum}
\graphicspath{ {figures/} }
\usepackage[font=footnotesize,labelfont=bf]{caption}

\hypersetup{
    colorlinks=true,
    linkcolor=blue,
}
\setlength{\parindent}{6ex}
\setlength{\parskip}{0.2ex}
\renewcommand{\baselinestretch}{1.2}

\author{Emre Geçit, Baran Yancı}
\begin{document}

    \title{\includegraphics[width=0.6\textwidth]{assets/favicon.png}\\ Software Architecture Description for \\  \textbf{Afet Bilgi}}
    \maketitle

    \newpage
    \makeatletter
	\renewcommand\tableofcontents{%
		\null\hfill\textbf{\Large\contentsname}\hfill\null\par
		\@mkboth{\MakeUppercase\contentsname}{\MakeUppercase\contentsname}%
		\@starttoc{toc}%
	}
	\makeatother

    \tableofcontents
    \listoffigures
    \listoftables
    \doublespacing

    \newpage

    \section{Introduction}
    \subsection{Purpose and Objectives}
    % TODO:
    The purpose of afetbilgi.com is to provide a centralized platform for users to access information about natural disasters and
    emergency response procedures. The website aims to be a reliable source of information that can help users prepare for,
    respond to, and recover from disasters.
    \subsection{Scope}
    % TODO:
    The scope of afetbilgi.com includes the following:
    \begin{itemize}
        \item Provide up-to-date information about natural disasters, including their causes, impacts, and possible response
        procedures.
        \item Provide important phone numbers, or the locations of important places or services to the people affected by the
        earthquake.
        \item Provide information about donation centers; such as blood, stem cells, money and so on to the people who want to
        help the people affected by the earthquake.
        \item Offer people to create help points and share the locations of the help points, with the help of the system.
    \end{itemize}

    The scope of afetbilgi.com does not include the following:
    \begin{itemize}
        \item Providing direct and instant communication services during a disaster event (this is the responsibility of
        official emergency response organizations).
        \item Offering financial assistance or physical aid during a disaster event (this is the responsibility of the official
        relief organizations).
    \end{itemize}
    \subsection{Stakeholders and Concerns}
    % TODO:
    The stakeholders of afetbilgi.com include the following:
    \begin{itemize}
        \item \textbf{Users:} They are concerned with accessing accurate information about natural disasters, preparing for
        emergencies, and communicating with others during a disaster event.
        \item \textbf{Emergency responders:} They are concerned with having access to timely and accurate information about
        disaster events to respond effectively.
        \item \textbf{Information providers:} People who voluntarily provide information about the earthquake. They can reach out
        to the maintainers of the system to provide information.
        \item \textbf{Maintainers / Developers:} People who maintain and develop the system. They develop the software and maintain
        it by verifying new information and adding it to the system.
    \end{itemize}

    \section{References}
    % TODO:
    The following references were consulted during the development of this SAD:
    \begin{itemize}
        \item ISO/IEC/IEEE 42010:2011, Systems and software engineering -- Architecture description
        \item Rozanski, N., \& Woods, E. (2011). Software systems architecture: working with stakeholders using viewpoints and 
        perspectives. Addison-Wesley Professional.
    \end{itemize}

    \section{Glossary}
    The following terms are used throughout this SAD and have the following meanings:
    \begin{itemize}
        \item \textbf{Disaster:} A sudden, catastrophic event that causes widespread damage or loss of life. Examples include 
        earthquakes, hurricanes, floods, and wildfires.
        \item \textbf{Emergency responder:} A person who is trained to respond to emergencies and provide assistance to 
        those in need. Examples include police officers, firefighters, and paramedics.
        \item \textbf{ISO/IEC/IEEE 42010:} A standard for software architecture description that provides a framework for 
        describing the architecture of a software-intensive system.
        \item \textbf{Relief organization:} An organization that provides aid and support to communities affected by disasters. 
        Examples include the Red Cross and UNICEF.
        \item \textbf{Scope:} The boundaries of a project or system, including what is included and what is excluded from
        consideration.
        \item \textbf{Stakeholder:} A person or group who has an interest in the success of a project or system. Examples include
        users, customers, and investors.
        \item \textbf{Viewpoint:} A perspective on the architecture of a system that emphasizes certain aspects or concerns.
        Examples include functional, data, and deployment viewpoints.
    \end{itemize}

    \section{Architectural Views}
    \subsection{Context View}
    \subsubsection{Stakeholders' Uses}
    There are four main stakeholders of the system: the users, emergency responders, information providers and the programmers.
    The users use this view to understand the general structure of their interaction with the system. The emergency responders
    use it to understand the general structure of the system to plan their response to the disaster. Information providers use it
    to supervise the processes their data goes through. Programmers use it to understand the general structure of the system to
    develop and maintain it.
    \subsubsection{Context Diagram}
    % TODO:
    \begin{figure}
    \includesvg[inkscapelatex=false, width=\textwidth ]{assets/ContextDiagram1-registered.svg}
    \caption{Context Diagram}
    \end{figure}
        As shown in the diagram, users can access the website through a web browser or mobile device. They can view information about
        natural disasters, emergency response procedures, and other related content. Administrators have additional privileges and can
        manage user accounts, update content on the website, and monitor site analytics.

        Third-party services are also used to provide additional functionality on the website. For example, a map API is used to
        display maps of affected areas and help users locate nearby emergency resources.

        Overall, the Context Diagram provides a high-level overview of how afetbilgi.com interacts with its external environment and
        stakeholders. This information is useful for understanding the scope of the system and identifying potential risks or
        dependencies that may impact its operation.

    \subsubsection{External Interfaces}
    \lipsum[1-1]
    \subsubsection{Interaction Scenarios}
    \lipsum[1-1]
    \subsection{Functional View}
    \subsubsection{Stakeholders' Uses}
    There are four main stakeholders of the system: the users, emergency responders, information providers and the programmers.
    The users use this view to understand the capabilities and limitations of the system. The emergency responders use it to
    determine how the system can be used to respond to the disaster. Information providers use it to understand what procedures
    are used and how they operate on their data. Programmers use it to understand the capabilities and limitations of the system
    to develop and maintain it.
    \subsubsection{Component Diagram}
    \subsubsection{Internal Interfaces}
    \lipsum[1-1]
    \subsubsection{Interaction Patterns}
    \lipsum[1-1]
    \subsection{Information View}
    \subsubsection{Stakeholders' Uses}
    There are four main stakeholders of the system: the users, emergency responders, information providers and the programmers.
    The users use this view to understand the data that is stored in the system. The emergency responders use it to understand
    the data that is stored in the system to plan their response to the disaster. Information providers use it to understand
    what data is stored in the system and how it is used. Programmers use it to understand the data that is stored in the system
    to develop and maintain it.
    \subsubsection{Database Class Diagram}
    \subsubsection{Operations on Data}
    \lipsum[1-1]
    \subsection{Deployment View}
    \subsubsection{Stakeholders' Uses}
    \lipsum[1-1]
    \subsubsection{Deployment Diagram}
    \subsection{Design Rationale}
    \lipsum[1-1]

    \section{Architectural Views for Suggestions to Improve the Existing System}
    \subsection{Context View}
    \subsubsection{Stakeholders' Uses}
    \lipsum[1-1]
    \subsubsection{Context Diagram}
    \subsubsection{External Interfaces}
    \lipsum[1-1]
    \subsubsection{Interaction Scenarios}
    \lipsum[1-1]
    \subsection{Functional View}
    \subsubsection{Stakeholders' Uses}
    \lipsum[1-1]
    \subsubsection{Component Diagram}
    \subsubsection{Internal Interfaces}
    \lipsum[1-1]
    \subsubsection{Interaction Patterns}
    \lipsum[1-1]
    \subsection{Information View}
    \subsubsection{Stakeholders' Uses}
    \lipsum[1-1]
    \subsubsection{Database Class Diagram}
    \subsubsection{Operations on Data}
    \lipsum[1-1]
    \subsection{Deployment View}
    \subsubsection{Stakeholders' Uses}
    \lipsum[1-1]
    \subsubsection{Deployment Diagram}
    \subsection{Design Rationale}
    \lipsum[1-1]
    
\end{document}